% JCP-style manuscript draft.  This source compiles without external template
% files; for a formal Elsevier submission, replace the class below with
% \documentclass[preprint,12pt]{elsarticle} and use elsarticle-num.bst.
\documentclass[11pt]{article}

\usepackage[a4paper,margin=2.45cm]{geometry}
\usepackage{newtxtext,newtxmath}
\usepackage{microtype}
\usepackage{amsmath,mathtools,bm}
\usepackage{booktabs,longtable,array}
\usepackage{graphicx}
\usepackage{enumitem}
\usepackage[hidelinks]{hyperref}
\usepackage[numbers,sort&compress]{natbib}
\usepackage{caption}
\usepackage{fancyhdr}
\usepackage{lastpage}

\graphicspath{{figures/}{../ncfv_first_reconstruction/}}
\setlength{\parskip}{0.35em}
\setlength{\parindent}{1.2em}
\setlength{\headheight}{14pt}
\setlength{\emergencystretch}{2em}
\setlist{nosep,leftmargin=*}
\renewcommand{\arraystretch}{1.12}
\captionsetup{font=small,labelfont=bf}
\pagestyle{fancy}
\fancyhf{}
\lhead{JCP-style manuscript draft}
\rhead{Differential-integration NCFV}
\cfoot{\thepage\ / \pageref{LastPage}}

\newcommand{\U}{\bm{U}}
\newcommand{\F}{\bm{F}}
\newcommand{\grad}{\bm{\nabla}}
\newcommand{\dd}{\mathrm{d}}
\newcommand{\Order}{\mathcal{O}}

\title{A differential-integration node-centred finite-volume framework for efficient high-order reconstruction on unstructured meshes}
\author{Author names and affiliations to be inserted}
\date{}

\begin{document}
\maketitle

\begin{abstract}
High-order node-centred finite-volume methods incur substantial costs from quadratic coefficient storage and repeated nonlinear flux evaluation.  Building on the author's master's thesis, this article develops a differential-integration formulation for unstructured dual meshes.  A complete quadratic least-squares problem supplies second-order gradients, while exact simplex moments and total-differential identities replace Hessian contractions by precomputed gradient weights.  The resulting operators recover nodal point values from cell averages and integrate physical fluxes without production Gauss points.  A chain-rule error analysis establishes the required flux-gradient accuracy under smoothness and bounded-stencil assumptions.  A frozen macroface dissipation matrix reduces the Riemann work to one call per primal edge, and an edge-owner MPI protocol shares one conservative flux.  The thesis cost model predicts 39.8\% fewer floating-point storage entries and a 10.31-fold reduction in multiplications for a representative tetrahedral stencil.  Its reported one-dimensional and extruded-vortex tests show near-third-order convergence.  Independently available DNDSR paired runs yield marching-time ratios of 4.94--17.36 and aggregate proportional-set-size reductions of 15.6--66.5\%.  The repository's present component-wise vortex convergence sequence does not yet establish stable third-order accuracy.  The formulation therefore identifies both the measured efficiency gains and the remaining accuracy conditions associated with nonplanar dissipation, viscous closure and boundary integration.
\end{abstract}

\noindent\textbf{Keywords:} high-order finite volume; node-centred discretization; unstructured mesh; quadrature-free integration; $k$-exact reconstruction; MPI

\section{Introduction}\label{sec:intro}

High-order methods remain attractive for vortex-dominated and wave-propagation flows because reduced numerical dissipation can lower the resolution required for a prescribed error.  Their use on unstructured meshes is nevertheless limited by reconstruction, flux integration and data-management costs \cite{wang2007review}.  The classical $k$-exact finite-volume route reconstructs a polynomial from cell averages and evaluates interface fluxes by sufficiently accurate quadrature \cite{barth1990,ollivier2002}.  The reconstruction stencil, the number of face quadrature points and the cost of nonlinear Riemann evaluations all increase with order and dimension.

Several research directions address parts of this cost.  Quadrature-free ADER finite-volume schemes use analytic integration and a local space--time procedure \cite{dumbser2007}; spectral-volume formulations have reconstructed a polynomial flux and analytically integrated it over subfaces \cite{wang2002sv,harris2008}.  Compact and multiple-correction methods instead seek to control stencil growth or reuse local differential information \cite{pont2017,wang2017}.  Vertex-centred median-dual formulations are of particular interest because their edge loop can reduce the number of control volumes and flux interfaces, and their high-order $k$-exact variants have been incorporated in production solvers \cite{setzwein2021,setzwein2022}.  These approaches establish that high-order reconstruction, analytic integration and edge-based data layouts are complementary rather than competing design choices.

The immediate methodological basis is Ma's master's thesis on an efficient third-order vertex-centred finite-volume method \cite{ma2026thesis}, especially Chapters 3--5.  It combines gradient-only reconstruction, differential integration and a macroface dissipation approximation.  The present article develops that formulation in a unified notation, makes its consistency assumptions explicit, and relates its analytical cost model and reported tests to the independently available DNDSR implementation evidence.  The underlying gradient-only method is therefore attributed to the thesis.

Two antecedents clarify this construction.  Liu and Vinokur \cite{liu1998} derive polynomial moments on segments, triangles and tetrahedra, supplying the geometric integration identities.  Nishikawa and White \cite{nishikawa2023} eliminate stored second derivatives in a third-order cell-centred tetrahedral scheme using nodal gradients and projected derivatives.  Their solution unknowns are point values at cell centres.  Here the evolved unknowns are dual-cell averages associated with primal vertices: the anchor generally differs from the integration-simplex centroid, so first moments and the full quadratic moment contraction must be retained.  This difference explains the separate point-recovery operator and the general weights developed below.

The central observation is that a complete quadratic reconstruction determines a consistent first derivative, while its second-derivative coefficients need not be stored or evaluated at run time.  For each micro-simplex of a dual control volume or dual interface, the Hessian contribution can be written as a precomputed linear combination of gradients at its construction vertices.  The resulting formulation has the following components:

\begin{enumerate}
  \item a mixed-element node-centred dual geometry whose internal macrofaces correspond one-to-one with primal edges;
  \item a differential-integration identity that replaces quadratic coefficient storage and production Gauss-point storage with gradient weights;
  \item a macroface flux formulation and edge-owner MPI protocol that preserve a unique numerical flux across partition boundaries; and
  \item a paired comparison against a conventional quadratic-plus-quadrature path, reporting both performance/memory measurements and the present limits of the accuracy evidence.
\end{enumerate}

The aim is not to relabel every analytic integral as ``quadrature free.''  In this paper the term means that the \emph{production} differential path creates, stores and traverses no volume, internal-face or boundary-face quadrature points.  High-order quadrature is still used as an independent diagnostic and as the conventional baseline.  This distinction is essential when numerical accuracy and resource use are discussed together.

\section{Node-centred finite-volume formulation}\label{sec:method}

\subsection{Dual control volumes and semi-discrete system}

Let $\Omega_i$ be the dual control volume associated with primal node $i$, let $V_i=|\Omega_i|$, and let $e=(i,j)$ denote a primal edge.  The dual macroface $\Gamma_e$ separates the two endpoint control volumes.  In two dimensions, each cell--edge--node chain produces a micro-triangle formed by the node, edge midpoint and arithmetic mean of cell vertices.  In three dimensions, each cell--face--edge--node chain similarly produces a micro-tetrahedron using the node, edge midpoint, arithmetic face point and arithmetic cell point.  These are definitions of the dual geometry; for non-simplex primal cells the arithmetic points are not substituted by area or volume centroids.

For the inviscid conservation law $\partial_t\U+\grad\cdot\F(\U)=\bm{0}$, the node-centred semi-discrete system is
\begin{equation}
 \frac{\dd\overline{\U}_i}{\dd t}
 =-\frac{1}{V_i}\left(\sum_{e\in\mathcal{E}(i)}B_{ie}\widehat{\bm\Phi}_e
 +\widehat{\bm\Phi}^{\partial}_i\right),
 \label{eq:semidiscrete}
\end{equation}
where $\overline{\U}_i$ is a dual-cell mean, $B_{ie}\in\{-1,+1\}$ is the oriented node--edge incidence, $\widehat{\bm\Phi}_e$ is the area-integrated numerical flux, and $\widehat{\bm\Phi}^{\partial}_i$ is the boundary contribution.  The same topology is used by the differential and conventional paths; only the integration and storage representations differ.

\subsection{Quadratic reconstruction with gradient-only runtime storage}\label{sec:recon}

For a scalar conservative component, write the mean-preserving polynomial as
\begin{equation}
 p_i(\bm x)=\bar u_i+\sum_{\ell=1}^{N_q}a_{i\ell}\psi_{i\ell}(\bm x),
 \quad
 \psi_{i\ell}=b_\ell(\bm\xi)-\frac{1}{V_i}\int_{\Omega_i}b_\ell(\bm\xi)\,\dd V,
 \quad N_q=\frac{d(d+3)}{2}.
 \label{eq:zero-mean}
\end{equation}
Thus $\int_{\Omega_i}p_i\,\dd V=V_i\bar u_i$ for any coefficient vector.  This is the zero-mean construction of thesis Section 3.2 \cite{ma2026thesis}.  Its moment subtraction is essential because a dual-cell average generally differs from the value at its anchor vertex.

At each node, a weighted least-squares problem is constructed from neighbouring dual-cell means in a complete quadratic basis.  With local scaled coordinate $\bm\xi=\bm H_i^{-1}(\bm x-\bm x_i)$, $\bm H_i=\operatorname{diag}(\Delta x_i,\Delta y_i,\Delta z_i)$ contains the directional half-spans of the \emph{full dual control volume}.  The half-spans are accumulated from micro-simplex construction vertices; they are not taken from a primal-cell bounding box or a reconstruction-stencil bounding box.  On translated periodic boundaries, relative-coordinate lower/upper bounds are first united across periodic copies, which prevents the periodic-box width from contaminating a local scale.

Let $\bm b_{im}$ be the difference between the mean quadratic basis in neighbour control volume $m$ and the target-control-volume mean basis.  A weighted SVD supplies
\begin{equation}
 \bm a_i=\bm P_i
 \left[
  \overline{\U}_{m_1}-\overline{\U}_{i},\ldots,
  \overline{\U}_{m_s}-\overline{\U}_{i}
 \right]^{T},
\label{eq:quadratic-svd}
\end{equation}
More explicitly, with $(\bm A_i)_{m\ell}=V_m^{-1}\int_{\Omega_m}\psi_{i\ell}\,\dd V$ and positive diagonal weight matrix $\bm W_i$, the retained map is
\begin{equation}
 \bm P_i=(\bm W_i^{1/2}\bm A_i)^\dagger\bm W_i^{1/2},
 \qquad
 \bm P_i^{(g)}=[\,\bm I_d\ \bm 0\,]\bm P_i.
 \label{eq:retained-map}
\end{equation}
The thesis writes the unweighted full-rank map as $(\bm A_i^T\bm A_i)^{-1}\bm A_i^T$; \eqref{eq:retained-map} expresses the same least-squares selection with the weighting and SVD used by the implementation.  The quadratic columns participate before the derivative rows are selected.  Dropping those columns before inversion would change the least-squares problem.

The thesis begins with face-neighbour dual cells, equivalent to primal-edge neighbours, and expands through successive neighbour layers when necessary.  Its suggested three-dimensional stencil size of 14--18 for nine unknown coefficients is a practical starting point.  Stencil cardinality alone does not establish unisolvency: the rank and conditioning tests must also pass, particularly near walls and on anisotropic grids.
where $\bm P_i$ is accepted only if the complete quadratic system has full numerical rank and its condition number is below the specified threshold.  The nonconstant basis has five rows in two dimensions and nine rows in three dimensions; its pure-square entries retain the $\xi_\alpha^2/2$ convention.  Thus the least-squares solve is complete-quadratic in both modes.  The differential path stores only the first $d$ rows of $\bm P_i$, which yield the physical gradient
\begin{equation}
 \bm G_i=\grad\U_i=\bm H_i^{-1}\bm a_{1,i}.
 \label{eq:gradient}
\end{equation}
Thus, the method does not replace a quadratic reconstruction by a linear one; it removes only the unnecessary runtime representation of its second-derivative block.  The reference-length matrix is used consistently when evaluating the target and neighbour mean bases and when restoring physical gradients.  The scalar $V_i^{1/d}$ remains available for distance weights and CFL/wall scales but is not substituted for $\bm H_i$ in the polynomial basis.

The dual-cell average and point value are related by a precomputed stencil,
\begin{equation}
 \overline{\U}_i=\U_i+\sum_{m\in\mathcal{S}_i}\bm w_{im}\cdot\bm G_m,
 \qquad
 \U_i=\overline{\U}_i-\sum_{m\in\mathcal{S}_i}\bm w_{im}\cdot\bm G_m.
 \label{eq:point-recovery}
\end{equation}
Here and below a vector weight acts independently on every conservative component.  Equation \eqref{eq:point-recovery} is also the recovery operation used before constructing interface states.

\subsection{Offline compression of simplex moments}\label{sec:differential}

Consider a line, triangle or tetrahedral micro-simplex $K$ with $n$ construction vertices $\bm r_r$, expanded about anchor node $i$.  Define $\bm d_r=\bm r_r-\bm x_i$ and $\bm D=\sum_{r=1}^{n}\bm d_r$.  Every construction point is an affine combination of original primal nodes, so its gradient is interpolated by the same coefficients,
\begin{equation}
 \bm g(\bm r_r)=\sum_m C_{rm}\bm g_m.
\label{eq:affine-gradient}
\end{equation}
The exact first and second moments about that anchor are
\begin{equation}
 \int_K\bm d\,\dd\mu=\frac{|K|}{n}\bm D,\qquad
 \int_K\bm d\bm d^T\,\dd\mu=
 \frac{|K|}{n(n+1)}
 \left(\bm D\bm D^T+\sum_{r=1}^{n}\bm d_r\bm d_r^T\right).
 \label{eq:simplex-moments}
\end{equation}
Here $\dd\mu$ is the intrinsic measure of $K$, and $n=2,3,4$ for a segment, triangle and tetrahedron.  These are the moment identities used in thesis Eqs.~(3-12)--(3-15) \cite{ma2026thesis,liu1998}.  They retain $\bm D\bm D^T$ even when the anchor is outside a surface simplex.  Only a centroid-based expansion has $\bm D=\bm 0$.
For a quadratic component $u$, the total differential gives
\begin{equation}
 \bm{\mathcal H}^{(u)}_i\bm d_r=\bm g(\bm r_r)-\bm g_i+\Order(h^2),
 \label{eq:total-differential}
\end{equation}
where $\bm{\mathcal H}^{(u)}_i$ is the Hessian and $\bm g=\grad u$.  Substitution into the exact moments of a quadratic polynomial produces
\begin{equation}
 \int_K u\,\dd V=|K|u_i+\sum_m\bm\ell^{K}_{im}\cdot\bm g_m .
 \label{eq:volume-functional}
\end{equation}
In the implemented affine representation, the weight is explicitly
\begin{equation}
\begin{split}
 \bm\ell^{K}_{im}={}&\frac{|K|}{n}\bm D\,\delta_{im}\\
 &+\frac{|K|}{2n(n+1)}
 \left[
   \bm D\sum_{r=1}^{n}(C_{rm}-\delta_{im})
  +\sum_{r=1}^{n}\bm d_r(C_{rm}-\delta_{im})
 \right].
\end{split}
\label{eq:differential-weight}
\end{equation}
Equation \eqref{eq:differential-weight} is evaluated once during geometry initialization and its nonzero contributions are merged by support-node index.  It has the same form for a volume micro-simplex and a surface micro-simplex, with $|K|$ interpreted as the appropriate line, area or volume measure.  The control-volume weights, divided by $V_i$, give the $\bm w_{im}$ in \eqref{eq:point-recovery}; therefore point recovery is a direct compressed integral, not an extrapolation from sampled Gauss values.

No volume, internal-face or boundary-face Gauss point is created after this compression.  The conventional path, in contrast, retains all quadratic coefficients and evaluates them at fourth-order quadrature locations.  The distinction is material: the efficient method removes both the location arrays and the repeated polynomial/flux evaluations that a point loop would require.

\subsection{Polynomial exactness and local accuracy of point recovery}\label{sec:local-accuracy}

The mechanism of thesis Section 3.2.4 \cite{ma2026thesis} can be stated with explicit error bounds.  Suppose a shape-regular family of micro-simplices has diameter $O(h)$, the interpolation reproduces affine functions, and the stencil coefficients remain uniformly bounded.  For a quadratic $u$, its gradient is affine, so exact nodal gradients give
$\bm g(\bm r_r)-\bm g_i=\bm{\mathcal H}^{(u)}\bm d_r$ exactly.  Substitution into \eqref{eq:simplex-moments} proves quadratic exactness of \eqref{eq:volume-functional}.  This argument includes off-centre anchors.

For a smooth $u\in C^3$, define the differential functional
$Q_{K,i}(u_i,\{\bm g_m\})=|K|u_i+\sum_m\bm\ell^K_{im}\cdot\bm g_m$.
The Taylor remainder and the moment-weight scaling give
\begin{equation}
 \left|Q_{K,i}(u_i,\{\bm g_m\})-\int_Ku\,\dd\mu\right|
 \le C|K|h^3,\qquad
 \sum_m\|\bm\ell^K_{im}\|\le C|K|h.
 \label{eq:local-bound}
\end{equation}
Consequently, using numerical gradients $\bm g_{m,h}=\bm g_m+O(h^2)$ adds at most $O(|K|h^3)$ to the integral.  Affine interpolation at an edge midpoint, arithmetic face point or arithmetic cell point retains this $O(h^2)$ gradient accuracy, provided all participating distances are $O(h)$.

After summing micro-volumes and dividing by $V_i$, point recovery satisfies
\begin{equation}
 |u_{i,h}-u(\bm x_i)|
 \le |\bar u_{i,h}-\bar u_i|+Ch^3.
 \label{eq:recovery-error}
\end{equation}
Thus second-order gradients suffice for third-order point recovery from sufficiently accurate averages.  This is a local reconstruction statement: evolved cell averages carry their own spatial and temporal errors.  A face-integral error $O(h^{d+2})$, divided by a cell volume $O(h^d)$, only supplies an $O(h^2)$ residual bound without further cancellation.  Stability and inter-face error structure must therefore be addressed before inferring third-order solution convergence from polynomial exactness alone.

\subsection{Nonlinear flux mapping and chain-rule accuracy}\label{sec:chain-rule}

Thesis Section 4.1 \cite{ma2026thesis} provides the step linking reconstruction to physical-flux integration.  Let $\U_h-\U=O(h^p)$ and $\bm G_{\alpha,h}-\partial_\alpha\U=O(h^q)$, where $\bm G_{\alpha,h}$ denotes the component vector for derivative direction $\alpha$.  Let $\bm A_\beta(\U)=\partial F_\beta/\partial\U$ be bounded and Lipschitz on the range of states.  Then
\begin{equation}
\begin{split}
 \bm E_{\alpha\beta}={}&\bm A_\beta(\U_h)\bm G_{\alpha,h}
       -\partial_\alpha F_\beta(\U),\\
 \bm E_{\alpha\beta}={}&[\bm A_\beta(\U_h)-\bm A_\beta(\U)]\bm G_{\alpha,h}
       +\bm A_\beta(\U)(\bm G_{\alpha,h}-\partial_\alpha\U)
       =O(h^{\min(p,q)}).
\end{split}
\label{eq:chain-error}
\end{equation}
Also $F_\beta(\U_h)-F_\beta(\U)=O(h^p)$.  With $(p,q)=(3,2)$, the physical-flux value is third-order accurate and its chain-rule gradient is second-order accurate, precisely the inputs required by \eqref{eq:local-bound}.  For compressible Euler flow these statements apply to smooth states with density and pressure bounded away from zero.  They do not apply across a shock or after a limiter activates.  No nonzero derivative assumption is required for this upper bound; a vanishing Jacobian component may improve the local error.

\subsection{Macroface state, physical-flux and dissipation assembly}\label{sec:macroface}

Let $S_e$ be the scalar sum of micro-surface measures and $\bm A_e$ the oriented vector sum.  On a nonplanar macroface, $S_e$ and $\|\bm A_e\|$ are deliberately kept distinct.  For an interior macroface, the side-averaged states are
\begin{equation}
 \overline{\U}_{e,L}=\U_i+\frac{1}{S_e}\sum_m\bm\omega^{L}_{em}\cdot\bm G_m,
 \qquad
 \overline{\U}_{e,R}=\U_j+\frac{1}{S_e}\sum_m\bm\omega^{R}_{em}\cdot\bm G_m,
 \label{eq:face-states}
\end{equation}
where $\bm\omega_{em}^{L/R}$ are accumulated versions of \eqref{eq:differential-weight}.  The corresponding physical-flux integral for side $s\in\{L,R\}$ is
\begin{equation}
 \bm\Phi^c_{e,s}=
 \bm F(\U_{i_s})\cdot\bm A_e+
 \phi_{i_s}\sum_m\sum_{\alpha,\beta}
 W^{s}_{em,\alpha\beta}\,
 \partial_\alpha F_\beta(\U_m),
 \label{eq:physical-flux-functional}
\end{equation}
where $W^{s}_{em,\alpha\beta}$ is the precomputed flux-gradient weight, $\phi_{i_s}$ is the side limiter factor and $\partial_\alpha F_\beta(\U_m)$ is calculated from the already available state gradient.  The implementation caches these physical-flux gradients for owned and ghost support nodes once per residual evaluation; this replaces repeated chain-rule work in every edge loop.

The single Riemann call per macroface is used only for the nonlinear dissipative correction:
\begin{equation}
 \widehat{\bm\Phi}_e =
 \frac{1}{2}\left(\bm\Phi^c_{e,L}+\bm\Phi^c_{e,R}\right)
 +S_e\left[
 \bm F^{\mathrm{num}}(\overline{\U}_{e,L},\overline{\U}_{e,R};\bm n_e)
 -\frac{1}{2}\left(
 \bm F(\overline{\U}_{e,L})\cdot\bm n_e+
 \bm F(\overline{\U}_{e,R})\cdot\bm n_e\right)
 \right],
 \label{eq:macroface-flux}
\end{equation}
with $\bm n_e=\bm A_e/\|\bm A_e\|$.  Equation \eqref{eq:macroface-flux} prevents the central physical flux from being approximated by a single state evaluation while avoiding a Riemann solve at every micro-surface point.  Roe's approximate solver \cite{roe1981} is one available flux option; the actual flux selected in a verification case must always be reported.

For the Roe--Pike form used in thesis Section 4.2.1 \cite{ma2026thesis}, define $\bm B_e=|\widetilde{\bm A}_e|=\bm R_e|\bm\Lambda_e|\bm R_e^{-1}$, with any entropy correction included in $|\bm\Lambda_e|$.  The correction in \eqref{eq:macroface-flux} becomes
\begin{equation}
 \bm\Phi^d_e=-\frac{S_e}{2}\bm B_e
       (\overline{\U}_{e,R}-\overline{\U}_{e,L}).
 \label{eq:frozen-roe}
\end{equation}
The Roe velocity and enthalpy are formed from the two macroface means, e.g.
$\widetilde{\bm v}=(\sqrt{\rho_L}\bm v_L+\sqrt{\rho_R}\bm v_R)/
(\sqrt{\rho_L}+\sqrt{\rho_R})$, with the analogous formula for $\widetilde H$.
Freezing this matrix over the macroface permits it to be taken outside the jump integral.  This is the algebraic reason that one characteristic decomposition replaces the repeated decompositions at quadrature points.

The approximation deserves a geometry qualification.  On a planar smooth face, if $\U_R-\U_L=O(h^3)$ and $\|\bm B(\bm x)-\bm B_e\|=O(h)$, the change in the integrated dissipative term is $O(S_eh^4)$.  On a nonplanar macroface, microface normals may differ by $O(1)$ under refinement, so the second assumption does not follow from state smoothness.  The thesis's single effective normal is retained as the discrete model, while its order on general mixed grids is assessed by full-scheme tests.  This prevents the small-jump argument from being used as an unconditional global-order proof.

Each macroface owns one sorted, de-duplicated \texttt{efficientStencil}.  A stencil entry contains left/right state-gradient weights, left/right flux-gradient matrices and the scalar weight used for surface-gradient recovery.  Zero-order state and flux contributions are not copied into every entry: they are represented once through the two anchor nodes, $S_e$ and $\bm A_e$.  This structure makes the inner loop a contiguous traversal of synchronized point values, state gradients and cached physical-flux gradients rather than a search through several sparse arrays.

\subsection{Limiter, boundary treatment and residual dataflow}\label{sec:runtime}

For shock control, the efficient limiter constrains the macroface averages in \eqref{eq:face-states}, not a primal-edge midpoint value.  Bounds are constructed from endpoint point values $\U_i$ and $\U_j$, and a single factor $\phi_i$ multiplies the \emph{whole} gradient correction of side $i$.  Scaling every contributing support-node gradient by that node's factor would define a different interface mean and is not used.  Point recovery itself remains unlimited; the factor is applied to interface averages, physical-flux-gradient integrals and viscous surface gradients.

For a scalar component, write the unlimited mean correction as $\delta_{ie}=\bar u_{e,i}-u_i$, and let $u^{\min}_{ij}=\min(u_i,u_j)$ and $u^{\max}_{ij}=\max(u_i,u_j)$.  An explicit statement of the endpoint-bounded limiter discussed in thesis Section 3.3 \cite{ma2026thesis} is
\begin{equation}
 \phi_{ie}=
 \begin{cases}
  \min(1,(u^{\max}_{ij}-u_i)/\delta_{ie}),&\delta_{ie}>0,\\
  \min(1,(u^{\min}_{ij}-u_i)/\delta_{ie}),&\delta_{ie}<0,\\
  1,&\delta_{ie}=0,
 \end{cases}
 \quad \phi_i=\min_{e\in\mathcal E(i)}\phi_{ie}.
 \label{eq:limiter-explicit}
\end{equation}
The bounded mean is $u_i+\phi_i\delta_{ie}$; \eqref{eq:face-states} displays the unlimited case.  A system implementation must specify how component factors and positivity protection are combined.  Endpoint bounds can activate even in smooth non-monotone regions, so formal smooth-order arguments assume an inactive limiter.

Boundary micro-surfaces use a compact affine stencil with a scalar integration weight.  The inviscid boundary flux is evaluated at each support-node state by the boundary flux function and then integrated by the affine weights.  This ordering matters for nonlinear boundary conditions: interpolating states first and applying the boundary flux once would be a different discrete operator.  When viscous terms are enabled, the same boundary stencil yields the averaged interior gradient; the current inviscid efficiency claims do not imply a verified third-order viscous scheme.

One efficient residual evaluation executes the following dataflow:
\begin{enumerate}
 \item pull dual-cell means; apply the precomputed first-derivative rows of \eqref{eq:quadratic-svd}; pull gradients;
 \item recover nodal point values with \eqref{eq:point-recovery}; pull point values and, when enabled, limiter factors;
 \item compute physical-flux gradients once per local/ghost support node;
 \item on each owned edge, evaluate \eqref{eq:physical-flux-functional} on both sides, form \eqref{eq:macroface-flux}, and compute the edge spectral-radius contribution;
 \item pull ghost-edge fluxes and spectral radii; assemble the signed node--edge divergence in \eqref{eq:semidiscrete}; then apply the periodic quotient average where relevant.
\end{enumerate}
Only node means, gradients, point values, limiter factors and edge flux/radius fields participate in runtime communication.  Micro-geometry and Gauss point fields do not.

\subsection{Viscous extension and affine boundary integration}\label{sec:viscous-extension}

Thesis Section 4.2.2 \cite{ma2026thesis} proposes a face-averaged viscous closure.  With our convention $\bm G\in\mathbb R^{d\times n_{\mathrm{var}}}$, its gradient correction can be written
\begin{equation}
 \bm G_e^{\mathrm{corr}}=\frac{\overline{\bm G}_{e,L}+\overline{\bm G}_{e,R}}{2}
  +\frac{\bm n_e(\overline{\U}_{e,R}-\overline{\U}_{e,L})^T}{2\tilde h_e},
 \qquad \tilde h_e=\frac{\min(V_i,V_j)}{S_e}.
 \label{eq:viscous-jump}
\end{equation}
The normal-jump contribution couples the two states in addition to averaging gradients.  For primitive variables $\bm q=\bm q(\U)$, a chain rule maps $\bm G_e^{\mathrm{corr}}$ to $\bm G_{q,e}^{\mathrm{corr}}$, and the proposed integral has the form
\begin{equation}
 \bm\Phi^{v}_e\simeq
 S_e\bm F^v\!\left(
       \frac{\bar{\bm q}_{e,L}+\bar{\bm q}_{e,R}}2,\bm G_{q,e}^{\mathrm{corr}},\bm n_e\right).
 \label{eq:viscous-mean}
\end{equation}
These formulas document the thesis extension; the paired runtime results below concern the inviscid implementation.  A second-order gradient does not provide a general third-order viscous residual.  In addition, moving a nonlinear constitutive law outside an integral is an approximation: temperature-dependent transport coefficients and products of variable velocity and gradients require a separate consistency analysis.

For a planar boundary triangle $T$ with three construction vertices and an affine numerical-flux interpolant $f_b$, the correctly normalized integral is
\begin{equation}
 \int_T f_b\,\dd S=\frac{|T|}{3}\sum_{r=1}^3 f_b(\bm r_r),
 \qquad
 f_b(\bm r_r)\simeq\sum_m C_{rm}f_b(\U_m;\bm n_T).
 \label{eq:boundary-affine}
\end{equation}
This expands the boundary construction of thesis Section 4.3.2 \cite{ma2026thesis}.  The $1/3$ factor is required by constant preservation; summing nodal values with $|T|$ alone would triple a constant flux.  The nodal flux includes the boundary-state construction before affine interpolation.  Curved boundaries also require normal and geometric-approximation control.  Concentrating a lower-order error on the boundary is insufficient on its own to guarantee third-order global accuracy: the PDE, stability, boundary closure and selected norm all matter.

\begin{table}[t]
\centering
\caption{Runtime representation of the two implementations.  The reconstruction solve is quadratic in both paths; ``rows'' are nonconstant reconstruction rows per solution component.}
\label{tab:representation}
\small
\resizebox{\linewidth}{!}{%
\begin{tabular}{lccc}
\toprule
Property & Conventional quadrature & Differential integration & Consequence \\
\midrule
2-D/3-D stored reconstruction rows & 5 / 9 & 2 / 3 & smaller dynamic state \\
Volume and surface points & stored and traversed & 0 in production & no point arrays \\
State at an integration location & quadratic evaluation & gradient functional & no Hessian storage \\
Riemann calls on one macroface & per surface point & one & fewer nonlinear flux calls \\
Partitioned macroface & local point loop & unique edge owner & one shared flux \\
\bottomrule
\end{tabular}
}
\end{table}

\subsection{MPI ownership and conservation}\label{sec:mpi}

Each primal edge has one deterministic owner, selected from the partition ownership of its adjacent cells.  The owner evaluates $\widehat{\bm\Phi}_e$ once; ranks holding ghost endpoints pull that result.  Node owners then apply the signed incidence in \eqref{eq:semidiscrete}.  The design ensures that the two endpoints of a cross-partition edge use exactly the same flux rather than independently recomputing rounded variants.  Point values and gradients are synchronized by the existing node pull maps.  The implementation does not retain raw pointers into communication arrays across allocation changes.

\section{Cost model and amortization of geometry preprocessing}\label{sec:cost}

The storage and arithmetic model in thesis Section 5.1 \cite{ma2026thesis}, printed pp.~57--59, makes the sources of efficiency quantitative.  It assumes a pure tetrahedral primal grid, five conservative variables, an inactive limiter and inviscid fluxes.  Let $s$ be the mean least-squares stencil size, $N_f$ the mean number of adjacent dual cells, and $N_e$ the mean number of primal tetrahedra around one edge.  Each dual macroface then contains $2N_e$ micro-triangles.  The thesis baseline uses three quadrature points per micro-triangle.  Contributions stored or computed once on a shared face are charged one half to each endpoint.

Under that counting convention, the numbers of floating-point storage entries per dual cell are
\begin{equation}
 M_T=16N_fN_e+14s+99,\qquad
 M_E=6N_fN_e+20N_f+8s+44.
 \label{eq:memory-model}
\end{equation}
The corresponding numbers of multiplications per spatial residual evaluation are
\begin{equation}
 \begin{split}
 C_T&=621N_fN_e+45s,\\
 C_E&=22.5N_fN_e+168.5N_f+15s+84.
 \end{split}
 \label{eq:operation-model}
\end{equation}
These reproduce thesis Tables 5.1--5.3.  The reduction from $45s$ to $15s$ reflects retaining three instead of nine coefficients for five equations.  The conventional term proportional to $N_fN_e$ includes state evaluations, nonlinear fluxes and entropy corrections at microface quadrature points.  The efficient path substitutes cached nodal flux gradients and support-weight contractions, with dissipation computed once per macroface.  The thesis calls this a per-time-step count; here it is interpreted as the spatial work per residual evaluation, to make the extra evaluations of a multistage integrator explicit.

\begin{table}[htbp]
\centering
\caption{Thesis cost model evaluated at $s=15$, $N_f=12$, $N_e=5$.  Counts are per dual cell; the ratios are recalculated from thesis Table 5.3 \cite{ma2026thesis}.}
\label{tab:cost-model}
\begin{tabular}{lrrr}
\toprule
Quantity & Conventional & Differential & Interpretation\\
\midrule
Floating-point storage entries & 1269 & 764 & 39.80\% reduction\\
Multiplications per residual & 37935 & 3681 & $C_T/C_E=10.31$\\
\bottomrule
\end{tabular}
\end{table}

The model predicts arithmetic work and algorithmic arrays, rather than elapsed time or process memory.  It excludes additions, divisions, square roots, cache behaviour, index storage, communication buffers and allocator overhead.  Moreover, the DNDSR baseline below uses six surface quadrature points per micro-triangle and its current fused stencil layout, so the thesis counts cannot be inserted as measured DNDSR bytes or timings.

For a fixed mesh, let $P_E,P_T$ be preprocessing times and $c_E,c_T$ measured costs of one residual evaluation.  A run with $N_R$ residual evaluations satisfies the simple model
\begin{equation}
 T_E=P_E+N_Rc_E,\qquad T_T=P_T+N_Rc_T.
 \label{eq:amortization}
\end{equation}
If $P_E>P_T$ and $c_T>c_E$, the break-even count is
$N_R>(P_E-P_T)/(c_T-c_E)$.  If $P_E\le P_T$ and $c_E<c_T$, no positive break-even delay occurs in this model.  This extension of the thesis count explains why a marching-only speed ratio should be accompanied by preprocessing cost for short runs or moving meshes.

At fixed polynomial order and bounded stencil sizes, both paths retain linear asymptotic work and storage in the number of nodes and edges.  The gain is a reduction in constants.  Runtime communication can become a larger fraction of the faster edge loop, so a serial arithmetic reduction does not by itself predict the parallel speedup.

\section{Implementation and verification protocol}\label{sec:implementation}

The algorithms were implemented as two selectable modes in the DNDSR NCFV module: \texttt{Efficient\allowbreak Differential} and \texttt{Traditional\allowbreak Quadrature}.  They share mesh preprocessing, least-squares stencil construction, limiters, time integration, flux family and MPI ownership.  The conventional mode retains all quadratic coefficients and uses fourth-order rules (14 points per micro-tetrahedron and 6 per micro-triangle in the reported three-dimensional setting).  The differential mode retains the gradient block and geometry/flux weights only.

The verification evidence used here consists of: (i) micro-geometry closure, (ii) polynomial-reproduction audits with independently evaluated reference quadrature, (iii) owner/ghost and periodic consistency checks, (iv) static reconstruction on a translated isentropic vortex, and (v) completed, paired time-marching tests.  Production tests retain no differential-mode Gauss points; diagnostic quadrature is isolated from production storage.  Geometry closure errors are approximately $10^{-16}$ in the reported periodic meshes, and parallel reproductions agree to roundoff in the corresponding audits.

\begin{figure}[t]
\centering
\includegraphics[width=0.98\linewidth]{convergence.png}
\caption{Static first-reconstruction errors of the differential mode for a periodic isentropic vortex.  The left panel measures recovered nodal values and the right panel measures first derivatives.  This is a $t=0$ reconstruction check, not an end-to-end convergence demonstration.}
\label{fig:static-reconstruction}
\end{figure}

\section{Numerical evidence}\label{sec:results}

\subsection{Accuracy and application results reported in the thesis}\label{sec:thesis-results}

The thesis contains additional completed tests beyond the repository evidence available for the previous manuscript draft.  The values in Tables~\ref{tab:thesis-1d} and \ref{tab:thesis-vortex} are transcribed from its numbered Tables 5.5 and 5.6 (printed pp.~61--62; PDF pages 75--76) \cite{ma2026thesis}.  They are source-reported results, not newly executed simulations.  The nominal convergence rate is $p_{\mathrm{obs}}=\log(E_h/E_{h/2})/\log 2$; the rounded source errors reproduce the stated rates to the displayed precision.

\begin{table}[htbp]
\centering
\caption{One-dimensional $L_1$ errors reported in thesis Table 5.5 \cite{ma2026thesis}.  Advection--diffusion is advanced for 10.5 periods, and Burgers' equation to $t=0.1$.}
\label{tab:thesis-1d}
\begin{tabular}{rrrrr}
\toprule
$N$ & Advection--diffusion error & Rate & Burgers error & Rate\\
\midrule
50 & $2.94\times10^{-2}$ & -- & $7.41\times10^{-4}$ & --\\
100 & $3.73\times10^{-3}$ & 2.98 & $9.72\times10^{-5}$ & 2.93\\
200 & $4.58\times10^{-4}$ & 3.02 & $1.37\times10^{-5}$ & 2.83\\
400 & $5.70\times10^{-5}$ & 3.01 & $1.65\times10^{-6}$ & 3.06\\
800 & $7.11\times10^{-6}$ & 3.00 & $2.10\times10^{-7}$ & 2.97\\
\bottomrule
\end{tabular}
\end{table}

\begin{table}[htbp]
\centering
\caption{Extruded isentropic-vortex results reported in thesis Table 5.6 \cite{ma2026thesis}: periodic domain $[0,10]\times[0,10]\times[0,4]$, $t=2$.  The source labels its metric as an $L_1$ error, without specifying the component in the table heading.}
\label{tab:thesis-vortex}
\begin{tabular}{rrrr}
\toprule
Nominal $h$ & Control volumes & Reported $L_1$ error & Rate\\
\midrule
1 & 745 & $3.36\times10^{-3}$ & --\\
$1/2$ & 4752 & $3.89\times10^{-4}$ & 3.11\\
$1/4$ & 33694 & $4.74\times10^{-5}$ & 3.04\\
$1/8$ & 252780 & $5.64\times10^{-6}$ & 3.07\\
\bottomrule
\end{tabular}
\end{table}

Figure~\ref{fig:thesis-convergence} redraws these tabulated values.  The one-dimensional and extruded-vortex series support the thesis's reported near-third-order behaviour for those configurations.  The extruded vortex has no imposed spanwise variation, so it exercises a three-dimensional mesh and operator without testing a fully three-dimensional solution field.  Its $L_1$ metric and nominal mesh size differ from the component-wise $L_2$ metric and measured $h$ used by the current repository sequence.  Although control-volume counts coincide for the prism sequence, equality of counts does not establish equality of flux, time-step, initialization or error-evaluation settings.

\begin{figure}[htbp]
\centering
\includegraphics[width=\linewidth]{thesis_convergence.png}
\caption{Convergence redrawn from thesis Tables 5.5--5.6 \cite{ma2026thesis}.  Dashed lines are third-order guides anchored at the finest value of each dataset; they are not additional simulation results.  No repository error series is mixed into these source-reported curves.}
\label{fig:thesis-convergence}
\end{figure}

Thesis Sections 5.2.3--5.2.5 also report three application studies.  The two-dimensional Riemann problems use an extruded $400\times400\times5$ mesh, testing interacting shocks and contact discontinuities.  The laminar flat-plate test uses $M_\infty=0.2$, $Re_L=10^5$, $Pr=0.72$, an adiabatic wall, a $200\times100\times5$ grid and pseudo-time CFL 30; its velocity and skin-friction curves are compared with Blasius profiles.  These plots support the reported application behaviour, but neither supplies a viscous grid-convergence study.

For the CRM wing--body, the thesis reports approximately six million mesh entities and a 512-core calculation at $M_\infty=0.85$ and $\alpha=2^\circ$, with $C_L=0.4513$ and $C_D=0.0197$.  The reference values quoted by the thesis are 0.47 and 0.029.  Its narrative describes an inviscid calculation, also lists $Re=5\times10^6$ and an adiabatic wall, and attributes the drag deficit to omitted turbulence modelling.  Those settings require clarification before a controlled validation claim can be made.  We retain this as a source-reported large-mesh demonstration; it supplies neither new experimental data nor a measured parallel-scaling curve.

\subsection{Current repository accuracy evidence and its scope}\label{sec:accuracy}

Figure~\ref{fig:static-reconstruction} shows that recovered nodal values decrease faster than third order in the $L_2$ norm on the two finest static reconstruction levels, while first derivatives converge close to second order, as expected from a complete quadratic reconstruction.  This does not itself prove the order of the full transient scheme.  In particular, point recovery combines the average and gradient errors through \eqref{eq:point-recovery}.

Table~\ref{tab:vortex-accuracy} reports the available end-to-end periodic isentropic-vortex calculation.  Both modes use $\mathrm{CFL}=0.5$, SSPRK3, no limiter, no viscosity and the same physical initial state; each solution reaches $t=2$.  The efficient calculation retains no Gauss points.  The conventional calculation is more accurate on these four levels, but both series exhibit an initially low observed order and the efficient method reaches only 2.43 for density and 2.68 for pressure in the finest $L_2$ comparison.  These data establish decreasing error and a useful controlled comparison, but they do \emph{not} establish stable third-order convergence of the complete method.  A publication-quality order claim requires temporal refinement/extrapolation, additional sufficiently fine grids, and verification with the intended standard Roe rather than a similarly named scalar-dissipation option.

\begin{table}[t]
\centering
\caption{Completed periodic-vortex results at $t=2$ on prism meshes.  $E$ and $T$ denote differential and conventional quadrature modes, respectively; entries are dual-volume-weighted point-value $L_2$ errors.  The observed order in parentheses is for the differential path.}
\label{tab:vortex-accuracy}
\small
\begin{tabular}{lccccc}
\toprule
Mesh & $h$ & $\rho_E$ & $p_E$ & $\rho_T$ & $p_T$ \\
\midrule
iv10 & 0.9210 & $5.1562\times10^{-2}$ & $7.5891\times10^{-2}$ & $4.8963\times10^{-2}$ & $7.3065\times10^{-2}$ \\
iv20 & 0.4683 & $2.6700\times10^{-2}$ (0.97) & $4.3078\times10^{-2}$ (0.84) & $2.2378\times10^{-2}$ & $3.6804\times10^{-2}$ \\
iv40 & 0.2360 & $5.7776\times10^{-3}$ (2.23) & $9.8230\times10^{-3}$ (2.16) & $4.5929\times10^{-3}$ & $7.5905\times10^{-3}$ \\
iv80 & 0.1186 & $1.0810\times10^{-3}$ (2.43) & $1.5524\times10^{-3}$ (2.68) & $8.3044\times10^{-4}$ & $1.1824\times10^{-3}$ \\
\bottomrule
\end{tabular}
\end{table}

\subsection{Paired performance and memory measurements}\label{sec:performance}

Table~\ref{tab:performance} compares actual marches to $t=2$ for prism, hexahedral and tetrahedral mesh families.  Each row is paired: both algorithms use the same mesh, MPI rank count and one thread per rank.  The reported time is the marching phase only, excluding initialization and file output; PSS is the aggregate proportional set size at the final state.  Rank counts change across mesh levels and therefore this table is \emph{not} a strong-scaling study.  It is a per-row algorithm comparison.

The conventional-to-differential marching-time ratio is 4.94--17.36.  Final PSS reductions are 15.6--66.5\%, with the largest reductions on tetrahedral meshes, where the conventional micro-simplex quadrature representation is especially numerous.  These measurements are consistent with the representation changes in Table~\ref{tab:representation}; they should not be extrapolated to viscous, limited, GPU or differently partitioned runs without direct measurement.  Full hardware, compiler, allocator and affinity metadata should accompany any submission version of this table.

\begin{longtable}{llrrrrrr}
\caption{Paired performance and memory measurements for full marches to $t=2$.  $t_E$ and $t_T$ are differential and conventional marching times; PSS is aggregated across ranks.}\label{tab:performance}\\
\toprule
Family & Mesh & ranks & nodes & $t_E$ (s) & $t_T$ (s) & $t_T/t_E$ & PSS reduction \\
\midrule
\endfirsthead
\toprule
Family & Mesh & ranks & nodes & $t_E$ (s) & $t_T$ (s) & $t_T/t_E$ & PSS reduction \\
\midrule
\endhead
Prism6 & iv10 & 1 & 745 & 0.743 & 7.772 & 10.47 & 15.6\% \\
Prism6 & iv20 & 1 & 4,752 & 11.752 & 122.166 & 10.40 & 44.1\% \\
Prism6 & iv40 & 8 & 33,694 & 29.943 & 280.974 & 9.38 & 48.7\% \\
Prism6 & iv80 & 32 & 252,780 & 176.211 & 1,260.767 & 7.15 & 42.3\% \\
\addlinespace
Hex8 & iv10 & 1 & 1,331 & 1.422 & 10.703 & 7.53 & 16.7\% \\
Hex8 & iv20 & 1 & 9,261 & 23.578 & 190.532 & 8.08 & 41.2\% \\
Hex8 & iv40 & 8 & 68,921 & 56.384 & 384.079 & 6.81 & 42.9\% \\
Hex8 & iv80 & 32 & 531,441 & 352.152 & 1,738.351 & 4.94 & 34.7\% \\
\addlinespace
Tet4 & iv10 & 1 & 1,331 & 1.825 & 31.681 & 17.36 & 40.8\% \\
Tet4 & iv20 & 1 & 9,261 & 32.385 & 540.866 & 16.70 & 66.5\% \\
Tet4 & iv40 & 8 & 68,921 & 76.646 & 1,123.901 & 14.66 & 65.5\% \\
Tet4 & iv80 & 32 & 531,441 & 412.866 & 4,866.953 & 11.79 & 55.3\% \\
\bottomrule
\end{longtable}

\section{Limitations and required completion studies}\label{sec:limitations}

The resource measurements and the absence of production quadrature storage are verified properties of the present implementation.  The following items remain before this draft can support a broad journal claim:
\begin{enumerate}
  \item Reconcile the thesis Tables 5.5--5.6 with archived meshes, solver revision, flux, CFL/time step, error component and norm definition.  The thesis reports near-third-order results; the present repository sequence measures a different error and does not yet reproduce an equivalent convergence test.
  \item Demonstrate asymptotic solution accuracy with a time-refinement study and more refined meshes; report $L_1$, $L_2$ and $L_\infty$ norms separately.
  \item Repeat inviscid studies with the intended characteristic Roe flux and document all entropy fixes and dissipation choices.  The current option called \texttt{Roe\_M2} follows a scalar local Lax--Friedrichs-type dissipation path and must not be described as standard Roe.
  \item Verify viscous accuracy independently.  The present differential treatment of inviscid volume/surface functionals does not by itself prove a third-order Navier--Stokes discretization.
  \item Add shock, boundary, curved-geometry and weak/strong-scaling studies, with compiler, CPU/GPU, memory and MPI-affinity details sufficient for reproduction.
\end{enumerate}

These limitations are not cosmetic qualifications.  They distinguish the measured reduction of quadrature work from a claim about the asymptotic accuracy of every configuration.  This is also consistent with the broader high-order unstructured-grid literature, where boundary treatment, nonlinear dissipation and reconstruction stencil quality must be validated jointly with formal polynomial consistency \cite{wang2007review,chen2022,setzwein2021}.

\section{Conclusions}\label{sec:conclusion}

Ma's thesis supplies the mathematical and algorithmic basis for gradient-only reconstruction and differential integration \cite{ma2026thesis}.  The expanded derivation shows how zero-mean least squares, exact simplex moments and an $O(h^2)$ gradient approximation give a local $O(h^3)$ point-recovery error when cell averages are sufficiently accurate.  Smooth nonlinear flux maps preserve the required value and derivative orders.  A frozen macroface Roe matrix explains the reduced characteristic work, with separate qualifications for nonplanar geometry, limiting and viscous terms.

A node-centred differential-integration finite-volume framework has been formulated and implemented for unstructured mixed-element meshes.  A full quadratic least-squares reconstruction supplies gradients, while total-differential identities turn volume and macroface integrals into precomputed gradient weights.  The production method stores no Gauss points and no quadratic reconstruction coefficients, and it uses one Riemann solve per primal-edge-associated macroface.  An edge-owner MPI protocol makes that flux unique across partitions.

The conventional-to-differential paired measurements show 4.94--17.36$\times$ faster marching and 15.6--66.5\% lower aggregate PSS in the current inviscid periodic tests.  Those efficiency results are substantial and directly traceable to the data representation.  In contrast, the current $t=2$ convergence sequence shows decreasing errors without stable third-order asymptotics.  The appropriate conclusion is therefore that the framework is an efficient, well-audited candidate for a third-order scheme, not that all accuracy requirements have already been met.  The completion studies in Section~\ref{sec:limitations} provide a concrete route from this manuscript draft to a submission-ready JCP article.

The thesis's separate $L_1$ datasets report rates near three, including 3.11, 3.04 and 3.07 for the extruded vortex.  Its analytical count predicts 39.8\% lower storage and a multiplication ratio of 10.31 for representative tetrahedral connectivity.  These source results strengthen the method's motivation and evidence base, while reproduction under matched settings remains necessary to explain the difference from the current repository convergence sequence.

\section*{Data availability}
The thesis source is Runzhi Ma, \emph{Research on the Efficient Third-order Vertex-centered Finite Volume Method}, Tsinghua University, April 2026, supplied by the author as a 94-page PDF.  Its extracted numerical values are distributed here as \texttt{thesis\_reported\_convergence.csv}; the supplemental figure is a redraw from those values.  Original solver outputs for the thesis tests were not supplied with the PDF.  The accompanying source map records printed-page and PDF-page locations.
The algorithm implementation, case configurations, diagnostics and numerical summaries used in this draft are maintained in the DNDSR repository.  The paired performance CSV is located at the repository root; the accuracy summaries are under \texttt{docs/reports/ncfv\_t2\_cfl05\_thesis\_20260907} and \texttt{docs/reports/ncfv\_traditional\_t2\_cfl05\_20260908}.  A public archival identifier and immutable revision should be added before submission.

\section*{Declaration of competing interest}
The authors declare no competing financial interests or personal relationships that could have appeared to influence the work reported in this paper.  Replace this placeholder with the authors' approved declaration before submission.

\section*{CRediT authorship contribution statement}
To be completed by the authors before submission.

\bibliographystyle{unsrtnat}
\bibliography{references}

\end{document}
